\section{Finite hidden-state bottlenecks give polynomial depth bounds}
\label{sec:finite-state-bottleneck}

Boolean inputs permit a positive result that has no analogue for a
continuously varying scalar input. If a short initial part of the teacher
produces only polynomially many distinct hidden states, its remaining
computation can be replaced by a lookup circuit. The remaining teacher
layers may be arbitrarily expansive. No additional condition is imposed
on the learner's prior.

\begin{proposition}[A finite first-layer state set]
\label{prop:first-layer-finite-state}
Suppose a tanh network in the present architecture computes
$f:\{-1,1\}^d\to\{-1,1\}$ and its first hidden vector has $C$
coordinates and at most $K$ distinct values on the cube. Then $f$ has a
depth-three threshold circuit with $O(KC(d+1))$ wires. Consequently it
has a globally unit-margin tanh realization of depth six, with width
and individual weights polynomial in $d,C,K$.
There is no bound on the depth, coefficients, or numerical stability
of the specified teacher's remaining computation.
\end{proposition}

\begin{proof}
For every attained hidden vector $v$, its label is well-defined because
all subsequent computation depends only on that vector. Tanh is
injective. Thus the condition that its $j$th coordinate equals $v_j$
can be tested by requiring the original affine preactivation to lie
between two thresholds that isolate its corresponding value from all
other attained preactivation values. There are only finitely many such
values. Choose separating endpoints strictly between consecutive values;
at an extreme, omit the unnecessary test or use a constant predicate.

Equality to $v$ is therefore a conjunction of at most $2C$ affine
threshold predicates of the original input. Use one threshold layer for
these predicates, a second for each state conjunction, and a final OR
over the positively labeled states. There are at most $K$ states and
each first-layer predicate uses at most $d+1$ wires. Constant functions
are immediate. Lemma~\ref{lem:threshold-circuit-tanh} gives the stated
bounded-weight tanh realization. In particular, arbitrarily small gaps
between attained real values cause no coefficient-size loss in the
resulting realization.
\end{proof}

The proposition already gives a sample polynomial in $d,C,K,L$ and the
statistical parameters at every learner depth $L\ge6$, simultaneously
for all $N\ge C_L(f)$ and all
$0<\sigma^2\le N^{-(1+\zeta)}$, for each fixed $\zeta>0$.
The next result allows a bottleneck after any fixed number of teacher
layers. Its proof uses binary precision, rather than a lower bound on
the gaps between all attained states.

\begin{lemma}[Fixed-depth tanh evaluation by threshold circuits]
\label{lem:fixed-prefix-tc-evaluation}
Fix $j\ge1$. Consider the first $j$ hidden layers of a tanh network
with widths at most $C$, input $(1,x)$ for $x\in\{-1,1\}^d$, and
weights of magnitude at most $B\ge1$. For every integer $p\ge1$, a
threshold circuit of depth depending only on $j$ and wire size
polynomial in $d,C,B,p$ produces a dyadic vector $\widehat h(x)$ such
that
\[
 \sup_x\|\widehat h(x)-h^j(x)\|_\infty\le2^{-p}.
\]
The polynomial degree may depend on $j$. Its advice consists of
polynomially many bits of the specified real parameters; no
computability assumption on those real parameters is required.
\end{lemma}

\begin{proof}
We use the theorem of Hesse--Allender--Barrington
\cite{hesse2002}: division and iterated multiplication of polynomially
many polynomial-bit integers have polynomial-size constant-depth
threshold circuits. Addition, iterated addition, and comparison also
have such circuits. Only their nonuniform consequence is needed here.

First consider a dyadic scalar $t=a/2^b$ with $|t|\le R$, $R\ge1$.
Tanh can be approximated to error $2^{-p}$ by a constant-depth
threshold circuit of size polynomial in $b,R,p$. Indeed, take the
degree-$m$ Taylor polynomial of $e^{2t}$, with
$m=\lceil40(R+p+2)\rceil$. Its uniform remainder is at most
\[
 e^{2R}\frac{(2R)^{m+1}}{(m+1)!}
 \le e^{2R}\left(\frac{2eR}{m+1}\right)^{m+1}
 <2^{-p-4}.
\]
The polynomial is the exact rational number $U/V$, where
\[
 V=2^{bm}m!,\qquad
 U=\sum_{k=0}^m (2a)^k2^{b(m-k)}\frac{m!}{k!}.
\]
All products and sums involve polynomially many integers with
$O(m(b+\log(R+1)+\log(m+1)))$ bits. Negative summands are handled
exactly; cancellation does not introduce numerical error.
The quoted arithmetic theorem computes these quantities at constant
depth. Since $U/V$ is within $2^{-p-4}$ of a nonnegative number,
$U/V+1\ge3/4$. Hence
\[
 \frac{U-V}{U+V}
\]
approximates $(e^{2t}-1)/(e^{2t}+1)=\tanh t$ to error less than
$2^{-p-2}$. A final rational division and dyadic rounding give the
required error. Clamping the result to $[-1,1]$ can only decrease its
distance from $\tanh t$.

To evaluate the prefix, round each parameter to absolute error
$2^{-q}$ and approximate every tanh to the same error, where
\[
 q=p+\left\lceil\log_2\bigl(4j(d+C+2)(2CB)^j\bigr)\right\rceil.
\]
Clamp all approximate hidden coordinates to $[-1,1]$.
If $E_s$ is the maximum hidden-coordinate error after layer $s$, then
\[
 E_1\le(d+2)2^{-q},\qquad
 E_s\le CB E_{s-1}+(C+1)2^{-q}\quad(2\le s\le j).
\]
Therefore
\[
 E_j\le j(d+C+2)(2CB)^j2^{-q}\le2^{-p}.
\]
The rounded preactivations have magnitude at most
$(d+C+1)(B+1)$ and polynomial bit length. Every affine calculation
uses only a constant number of the arithmetic operations above, and
each tanh uses the preceding scalar construction. Composing the fixed
number $j$ of stages preserves constant threshold depth and polynomial
wire size. The rounded matrices have polynomial total description
length. Their bits are fixed circuit advice, rather than quantities
the circuit is required to extract from arbitrary real numbers.
\end{proof}

\begin{theorem}[A finite bottleneck after a fixed prefix]
\label{thm:fixed-prefix-finite-state}
Fix $j\ge1$ and $\zeta>0$. Suppose a width-$C$, depth-$D>j$ tanh
teacher has weights bounded by $B\ge1$, signed output at least $1/2$
on the Boolean cube, and at most $K$ distinct vectors after its $j$th
hidden layer. There are a constant $q_j$ and a polynomial $p_j$, with
degrees depending only on $j$, such that $f$ has a globally
unit-margin tanh realization of depth $q_j$, width and weights at most
$p_j(d,C,B,K,D)$.

Consequently the original unrestricted Gaussian posterior has sample
complexity polynomial in
\[
 d,C,B,K,D,L,\varepsilon^{-1},\log(2/\delta)
\]
at every learner depth $L\ge q_j$, simultaneously for every
$N\ge C_L(f)$ and every $0<\sigma^2\le N^{-(1+\zeta)}$.
The same polynomial form holds at every learner depth $L\ge D$,
including the possible bounded-depth range below $q_j$.
The constants and polynomial degrees may depend on $j$ and $\zeta$,
but not on the remaining teacher depth $D-j$.
\end{theorem}

\begin{proof}
In the infinity norm, the teacher suffix from its $j$th hidden layer
to the linear readout is Lipschitz with constant at most
\[
 S=(CB)^{D-j}.
\]
Oppositely labeled inputs have outputs differing by at least one.
Consequently their $j$th hidden states are separated in infinity norm
by at least $S^{-1}$. Set
\[
 p=\lceil\log_2(32S)\rceil,\qquad \eta=2^{-p}.
\]
The required precision is only $O(1+D\log(2CB))$ bits, despite the
possibly exponentially small separation.

Apply Lemma~\ref{lem:fixed-prefix-tc-evaluation} to produce
$\widehat h(x)$ within $\eta$ of the true prefix state. For each
attained state $v$, hardcode a dyadic center $\widehat v$ within
$\eta$ of $v$, and hardcode its Boolean label. Output $+1$ exactly
when $\widehat h(x)$ lies within $3\eta$ in every coordinate of at
least one positive center. A positive input lies within $2\eta$ of
its own center. A negative input is at distance at least
$S^{-1}-2\eta\ge30\eta$ from every positive center. Thus the
lookup returns the correct label everywhere.

There are at most $K$ centers, each described by $O(Cp)$ bits.
Subtraction, comparison, a conjunction over coordinates, and a
disjunction over centers are constant-depth threshold operations of
polynomial size. Together with the fixed-depth prefix evaluation,
this gives a threshold circuit of depth depending only on $j$ and
size polynomial in $d,C,B,K,D$. Apply
Lemma~\ref{lem:threshold-circuit-tanh} and pad the resulting fixed
depth to a constant $q_j$. This proves the representation statement.

The fixed-intrinsic-depth learning theorem gives the claimed
all-admissible-width guarantee for $L\ge q_j$. To cover every
$L\ge D$, it remains only to consider $D<q_j$ and $L<q_j$.
There are finitely many such intrinsic teacher depths, bounded in terms
of $j$ alone. Double the supplied teacher's readout and apply the same
fixed-intrinsic-depth theorem directly to that teacher. Taking the
maximum of these finitely many polynomial bounds gives the stated
form. The prior throughout retains only its original global output
margin restriction and the observed training signs.
\end{proof}

In the representation proof it would suffice to cover the attained
prefix states by $K$ infinity-norm balls of radius $\eta$, with centers
chosen among those states. Use boxes of radius $4\eta$ in the lookup.
Opposite labels cannot occur in the same such ball, and the same
separation argument applies. Exact finite-state count is a particularly
simple sufficient condition; no separation of same-label states is
needed.

\begin{corollary}[Discrete input statistics]
\label{cor:discrete-statistic-bottleneck}
The following teachers satisfy Proposition~\ref{prop:first-layer-finite-state}.
\begin{enumerate}
\item If the nonconstant entries of first-layer row $j$ are
$A_{ji}/Q_j$, with integers $A_{ji}$ and $Q_j>0$, put
$R_j=\sum_{i=1}^d|A_{ji}|$. Its state count is at most
\[
 K\le\prod_{j=1}^C(2R_j+1).
\]
Its first-layer biases may be arbitrary real numbers. In particular,
first-layer input weights bounded by $B$ with common denominator $Q$
give $K\le(2dBQ+1)^C$. For fixed $C$ and polynomially bounded $B,Q$,
the sample bound is polynomial in input dimension and learner depth,
independently of all later teacher computation.
\item If the first-layer matrix has at most $k$ distinct nonconstant
input-column vectors, let $d_1,\ldots,d_k$ be their multiplicities.
Then
\[
 K\le\prod_{a=1}^k(d_a+1)\le(d+1)^k.
\]
For fixed $k$, the sample bound is polynomial in $d,C,L$ and the
statistical parameters, with no dependence on the precision or
magnitude of the specified teacher coefficients. Functions of a fixed
number of coordinate-block Hamming weights are examples.
\end{enumerate}
\end{corollary}

\begin{proof}
In the first case, the integer sum in row $j$ lies between $-R_j$
and $R_j$, so it has at most $2R_j+1$ values. A fixed bias and tanh
do not increase this count. In the second case, the first affine map
depends only on the sum of the Boolean inputs in each column group.
A group of size $d_a$ has exactly $d_a+1$ possible sums. All later
computation factors through the resulting state vector. Apply the
proposition and the fixed-intrinsic-depth learning theorem.
\end{proof}

These results use finite information passed by a teacher prefix,
rather than stability of the prior's networks. They do not imply a
joint polynomial theorem for an arbitrary growing prefix depth or
unrestricted state count. Even one bounded real affine projection can
take $2^d$ distinct values on the cube. Rounding that projection while
preserving a long suffix's output signs may require exponentially fine
precision in depth, so it does not establish the required polynomial
state bound for a general teacher.
